% !TEX root = MaCFP_2021_Report_Part_I_Experimental.tex
\chapter{Nomenclature}
\label{nomenclature}
\vspace*{-10mm}
\begin{tabbing}
	CAPA		\hspace{1in}\= controlled atmosphere cone calorimetry\\
	DSC						\>Differential scanning calorimetry\\
	FPA		        		\> Fire propagation apparatus\\
	MCC		        		\> Mircroscale combustion calorimetry\\
	TGA		        		\> Thermogravimetric analysis\\
	TM-DSC					\>Temperature-Modulated Differential scanning calorimetry\\
	TPS		        		\> Transient plane source\\
	
	\hspace{0.1in}            		\> \\
	{\bf  Parameters}       	\> \\
	$A$		        		\> Initial sample surface area\\
	$h_{\rm r}$        		\> Heat of decomposition reaction\\
	$\Delta H_{\rm c}$		\> Heat of combustion (energy released per gram of gaseous volatiles produced)\\
	$m$		        		\> Mass\\
	$\dot{m}''$				\> Mass loss rate (normalized by initial sample-area)\\
	$n$						\> Interval, defines number of points across which  average or standard deviation are calculated\\
	$N$		        		\> Number (of measurements, tests, or test locations)\\
	$\dot{q}$		     	\> Heat flow\\
	$q$		        		\> Integral heat flow\\
	$\dot{Q}''$				\> Heat release rate (normalized by initial sample-area)\\
	$t$						\> Time\\
	$T$						\> Temperature \\
	$X$               		\> Volume fraction \\
	
	\hspace{0.1in}            	\> \\
	{\bf Greek Letters}       	\> \\
	$\beta$               	\> Heating rate\\
	$\mu_{\rm char}$		\> Char yield\\
	$\sigma_{\rm mean}$   	 	\> Standard deviation of the mean \\
	
	\hspace{0.1in}            	\> \\
	{\bf Subscripts}      	    \> \\
	\hspace{0.1in}            	\> \\
	0     	                    \> Initial value \\
	$i$                         \> Time or temperature step\\
	$j$                         \> Test replicate\\
	$k$	                        \> Location replicate (e.g., multiple thermocouples used in same experiment)\\
	O$_2$                       \> Oxygen \\
	back         				\> Back surface\\
	endset         				\> Endset (of reaction)\\
	front         				\> Front surface\\
	onset                  		\> Onset (of reaction)\\
	tot                  		\> Total\\

	
\end{tabbing}




\chapter{Supplemental Information: All Measurement Data from Individual Labs}
\label{SI}
This Supplemental Information section has been prepared on behalf of the MaCFP Condensed Phase Working Group. It is written to provide supporting information to the main text of the predecisional draft report titled, `Preliminary Summary of Experimental Measurements'. All figures in this document have been prepared as described in the main text of this predecisional draft report. These figures have been prepared for subject matter experts to provide critical review and to ensure the integrity of the measurement data and related analysis submitted to the 2026 MaCFP-4 Condensed Phase Workshop. 

When error bars are provided, they are calculated as two standard deviations of the mean. This represents an estimate (95~\% confidence interval) of where the mean value is expected, given measurement data from replicate tests conducted under nominally the same conditions.

\newpage
\section{Thermogravimetric Analysis (TGA)}\label{App:TGA}
\subsection{Nitrogen Environment}
\label{TGA_N2}

\subsubsection{Heating Rates}

\input{../SCRIPT_FIGURES/TGA/TGA_Appendix_Heating_Rates.tex}

\newpage

\subsubsection{Individual Laboratory Results}

\input{../SCRIPT_FIGURES/TGA/TGA_Appendix_Individual_N2.tex}
\newpage

\subsubsection{Mass and Mass Loss Rate Curves}

\input{../SCRIPT_FIGURES/TGA/TGA_Appendix_Mass_MLR_N2.tex}
\newpage

\subsubsection{Peak and Onset Temperatures}

\noindent\textbf{$\beta$ = 5~K/min}

\begin{figure}[H]
    \centering
    \includegraphics[width=0.48\linewidth]{../SCRIPT_FIGURES/TGA/Tpeak_Average_N2_5K_MLR.pdf}
    \hfill
    \includegraphics[width=0.48\linewidth]{../SCRIPT_FIGURES/TGA/Tonset_Average_N2_5K.pdf}
    \caption{Peak normalized mass loss rate and onset temperature results for TGA tests conducted in pure nitrogen at a nominal heating rate of $\beta=5$~K/min.}
\end{figure}

\noindent\textbf{$\beta$ = 20~K/min}

\begin{figure}[H]
    \centering
    \includegraphics[width=0.48\linewidth]{../SCRIPT_FIGURES/TGA/Tpeak_Average_N2_20K_MLR.pdf}
    \hfill
    \includegraphics[width=0.48\linewidth]{../SCRIPT_FIGURES/TGA/Tonset_Average_N2_20K.pdf}
    \caption{Peak normalized mass loss rate and onset temperature results for TGA tests conducted in pure nitrogen at a nominal heating rate of $\beta=20$~K/min.}
\end{figure}

\newpage

\subsection{Oxidative Environments}
\label{TGA_O2}

\subsubsection{Individual Laboratory Results}

\input{../SCRIPT_FIGURES/TGA/TGA_Appendix_Individual_O2.tex}

\newpage

\subsubsection{Mass and Mass Loss Rate Curves}

\input{../SCRIPT_FIGURES/TGA/TGA_Appendix_Mass_MLR_O2.tex}

\newpage


\section{Differential Scanning Calorimetry (DSC)}
\subsection{Nitrogen Environment}
\label{DSC_N2}

\subsubsection{Average Heat Flow and Integral Heat Flow by Institution}

\input{../SCRIPT_FIGURES/DSC/DSC_Appendix_N2_Averages.tex}

\subsubsection{Individual Heat-of-Decomposition Analysis}

\input{../SCRIPT_FIGURES/DSC/DSC_Appendix_N2_Individual.tex}


\subsubsection{Inter-Laboratory Comparison}

\noindent\textbf{$\beta$ = 5~K/min}\par\vspace{0.5em}

\noindent
\begin{minipage}{0.49\textwidth}
    \centering
    \includegraphics[width=\linewidth]{../SCRIPT_FIGURES/DSC/DSC_Wood_N2_5K_HF.pdf}
\end{minipage}
\hfill
\begin{minipage}{0.49\textwidth}
    \centering
    \includegraphics[width=\linewidth]{../SCRIPT_FIGURES/DSC/DSC_Wood_N2_5K_HF_avg.pdf}
\end{minipage}

\par\vspace{0.5em}

\noindent
\begin{minipage}{0.49\textwidth}
    \centering
    \includegraphics[width=\linewidth]{../SCRIPT_FIGURES/DSC/DSC_Wood_N2_5K_iHF.pdf}
\end{minipage}
\hfill
\begin{minipage}{0.49\textwidth}
    \centering
    \includegraphics[width=\linewidth]{../SCRIPT_FIGURES/DSC/DSC_Wood_N2_5K_iHF_avg.pdf}
\end{minipage}

\par\vspace{1em}

\noindent\textbf{$\beta$ = 10~K/min}\par\vspace{0.5em}

\noindent
\begin{minipage}{0.49\textwidth}
    \centering
    \includegraphics[width=\linewidth]{../SCRIPT_FIGURES/DSC/DSC_Wood_N2_10K_HF.pdf}
\end{minipage}
\hfill
\begin{minipage}{0.49\textwidth}
    \centering
    \includegraphics[width=\linewidth]{../SCRIPT_FIGURES/DSC/DSC_Wood_N2_10K_HF_avg.pdf}
\end{minipage}

\par\vspace{0.5em}

\noindent
\begin{minipage}{0.49\textwidth}
    \centering
    \includegraphics[width=\linewidth]{../SCRIPT_FIGURES/DSC/DSC_Wood_N2_10K_iHF.pdf}
\end{minipage}
\hfill
\begin{minipage}{0.49\textwidth}
    \centering
    \includegraphics[width=\linewidth]{../SCRIPT_FIGURES/DSC/DSC_Wood_N2_10K_iHF_avg.pdf}
\end{minipage}

\par\vspace{1em}

\noindent\textbf{$\beta$ = 20~K/min}\par\vspace{0.5em}

\noindent
\begin{minipage}{0.49\textwidth}
    \centering
    \includegraphics[width=\linewidth]{../SCRIPT_FIGURES/DSC/DSC_Wood_N2_20K_HF.pdf}
\end{minipage}
\hfill
\begin{minipage}{0.49\textwidth}
    \centering
    \includegraphics[width=\linewidth]{../SCRIPT_FIGURES/DSC/DSC_Wood_N2_20K_HF_avg.pdf}
\end{minipage}

\par\vspace{0.5em}

\noindent
\begin{minipage}{0.49\textwidth}
    \centering
    \includegraphics[width=\linewidth]{../SCRIPT_FIGURES/DSC/DSC_Wood_N2_20K_iHF.pdf}
\end{minipage}
\hfill
\begin{minipage}{0.49\textwidth}
    \centering
    \includegraphics[width=\linewidth]{../SCRIPT_FIGURES/DSC/DSC_Wood_N2_20K_iHF_avg.pdf}
\end{minipage}


\subsection{Oxidative Environments}
\label{DSC_O2}

\subsubsection{Average Heat Flow and Integral Heat Flow}

\input{../SCRIPT_FIGURES/DSC/DSC_Appendix_O2_Averages.tex}


\newpage
\section{Micro-Scale Combustion Calorimeter (MCC)}
\subsection{Nitrogen Environment}
\label{MCC_N2}
\subsubsection{$\beta$ = 30~K/min}
\begin{minipage}{0.5\textwidth}
\begin{figure}[H]
	{\includegraphics[width=3.25in]{../SCRIPT_FIGURES/MCC/Average/FZJ_Wood_MCC_N2_30K.pdf}}\\
\end{figure}
\end{minipage} 
\begin{minipage}{0.35\textwidth}
\begin{figure}[H]
	{\includegraphics[width=3.25in]{../SCRIPT_FIGURES/MCC/Average/TUBS_Wood_MCC_N2_30K.pdf}}\\
\end{figure}
\end{minipage}

\subsubsection{$\beta$ = 45~K/min}
\begin{minipage}{0.5\textwidth}
\begin{figure}[H]
	{\includegraphics[width=3.25in]{../SCRIPT_FIGURES/MCC/Average/FZJ_Wood_MCC_N2_45K.pdf}}\\
\end{figure}
\end{minipage} 

\newpage
\subsubsection{$\beta$ = 60~K/min}
\begin{minipage}{0.5\textwidth}
\begin{figure}[H]
	{\includegraphics[width=3.25in]{../SCRIPT_FIGURES/MCC/Average/FZJ_Wood_MCC_N2_60K.pdf}}\\
\end{figure}
\end{minipage} 
\begin{minipage}{0.35\textwidth}
\begin{figure}[H]
	{\includegraphics[width=3.25in]{../SCRIPT_FIGURES/MCC/Average/IMT_Wood_MCC_N2_60K.pdf}}\\
\end{figure}
\end{minipage}\\

\begin{minipage}{0.5\textwidth}
\begin{figure}[H]
{\includegraphics[width=3.25in]{../SCRIPT_FIGURES/MCC/Average/Lund_Wood_MCC_N2_60K.pdf}}\\
\end{figure}
\end{minipage} 
\begin{minipage}{0.35\textwidth}
\begin{figure}[H]
{\includegraphics[width=3.25in]{../SCRIPT_FIGURES/MCC/Average/NIST_Wood_MCC_N2_60K.pdf}}\\
\end{figure}
\end{minipage}\\

\begin{minipage}{0.5\textwidth}
	\begin{figure}[H]
		{\includegraphics[width=3.25in]{../SCRIPT_FIGURES/MCC/Average/UDRI_Wood_MCC_N2_60K.pdf}}\\
	\end{figure}
\end{minipage} \\

\newpage
\subsection{Oxygen Environment}
\label{MCC_O2}
\subsubsection{$X_{\rm O_2}=0.02$ $\beta$ = 60~K/min}
\begin{minipage}{0.5\textwidth}
\begin{figure}[H]
{\includegraphics[width=3.25in]{../SCRIPT_FIGURES/MCC/Average/IMT_Wood_MCC_O2-2_60K.pdf}}\\
\end{figure}
\end{minipage} 
\vfill


\subsubsection{$X_{\rm O_2}=0.05$ $\beta$ = 60~K/min}
\begin{minipage}{0.5\textwidth}
\begin{figure}[H]
{\includegraphics[width=3.25in]{../SCRIPT_FIGURES/MCC/Average/IMT_Wood_MCC_O2-5_60K.pdf}}\\
\end{figure}
\end{minipage} 
\vfill


\subsubsection{$X_{\rm O_2}=0.10$ $\beta$ = 60~K/min}
\begin{minipage}{0.5\textwidth}
\begin{figure}[H]
{\includegraphics[width=3.25in]{../SCRIPT_FIGURES/MCC/Average/IMT_Wood_MCC_O2-10_60K.pdf}}\\
\end{figure}
\end{minipage} 
\vfill

\subsubsection{$X_{\rm O_2}=0.20$ $\beta$ = 60~K/min}
\begin{minipage}{0.5\textwidth}
\begin{figure}[H]
{\includegraphics[width=3.25in]{../SCRIPT_FIGURES/MCC/Average/IMT_Wood_MCC_O2-20_60K.pdf}}\\
\end{figure}
\end{minipage} 
\begin{minipage}{0.35\textwidth}
\begin{figure}[H]
{\includegraphics[width=3.25in]{../SCRIPT_FIGURES/MCC/Average/Lund_Wood_MCC_O2-20_60K.pdf}}\\
\end{figure}
\end{minipage}
\vfill
\begin{minipage}{0.5\textwidth}
\begin{figure}[H]
{\includegraphics[width=3.25in]{../SCRIPT_FIGURES/MCC/Average/NIST_Wood_MCC_O2-20_60K.pdf}}\\
\end{figure}
\end{minipage} \\


\subsubsection{Wood-char $X_{\rm O_2}=0.20$ $\beta$ = 20~K/min}
\begin{minipage}{0.5\textwidth}
	\begin{figure}[H]
		{\includegraphics[width=3.25in]{../SCRIPT_FIGURES/MCC/Average/Lund_Wood-char_MCC_O2-20_60K.pdf}}\\
	\end{figure}
\end{minipage} 
\begin{minipage}{0.35\textwidth}
	\begin{figure}[H]
		{\includegraphics[width=3.25in]{../SCRIPT_FIGURES/MCC/Average/NIST_Wood-char_MCC_O2-20_60K.pdf}}\\
	\end{figure}
\end{minipage}\\

\subsection{Baseline correction algorithm}
This section describes the baseline algorithm used to correct the raw MCC data, submitted by participants~\footnote{The algorithm was tuned for the dataset analyzed in this study. Its performance has not been validated for MCC data from other materials or experiments conducted under different conditions.}. The python scripts used to prepare this document, including the baseline correction, are available on the MaCFP repository. Figure~\ref{Fig:MCC_SI_baseline} shows the baseline algorithm on one of the submitted datasets, the algorithm is explained in more detail bellow. First, the data gets prepared in an initial step. Next, a region before the main HRR peak and one after gets identified, these regions are the starting points to search for the intervals used for the baseline correction. Finally a polynomial baseline correction is performed on the original data. 

 \begin{figure}[h]
	\centering
	\includegraphics[width=0.7\linewidth]{../SCRIPT_FIGURES/MCC/Baseline/FZJ_Wood_MCC-raw_N2_60K_R6.pdf}
	\caption{Baseline correction algorithm. The raw submitted data (light gray), gets prepared (dark gray). The algorithm then identifies the points immediately before and after the main HRR peak (green circles). These serve as a starting point to identify a region before (between the red stars) and after the peak (between the blue triangles) that is used for the baseline, the baseline is indicated in red and the corrected data is displayed in black.}
	\label{Fig:MCC_SI_baseline}
\end{figure}

The first step of the baseline algorithm is to prepare the data. Temperature data is converted to kelvin and all data is interpolated in 0.5~K intervals. Some of the raw data has an initial spike, which is removed by the algorithm, i.e.~in Fig.~\ref{Fig:MCC_SI_baseline} all data before the yellow cross gets ignored. Additionally, the HRR data is shifted (along the y-axis) such that all values are positive. Finally, a smoothed HRR column was generated by calculating the 20-point rolling mean of the HRR data. This column is referred to as the smoothed HRR in the following description.

In the second step regions before and after the main HRR peak are identified, i.e. the green dots in Fig.~\ref{Fig:MCC_SI_baseline} are determined. The point bebore the main peak is identified as the smallest index (above index 50) where the smoothed HRR exceeds 4\% of the peak HRR. If this index is smaller than 100, it is updated to the smallest index where the smoothed HRR exceeds 6\% of the peak HRR. If the coresponding index is still bellow 100, it is updated to 10\% of the peak HRR, at this point the identified index was always larger than 100. The region after the peak HRR is identified as the smallest index after the main peak, where the smoothed HRR is decreases bellow 9\% of the peak HRR. 

In the third step the data points used for the baseline correction located before the HRR peak are determined. This part of the algorithm searches for two distinct points, all the data between both points is included in the baseline fit. For the first point, i.e. the point closest to the start of the test, the algorithm identifies the earliest location at which the data are sufficiently flat. This is done by sliding an 80 point window between the start of the test and the green point found in step 2. For each window, the slope of a linear fit of the temperature-HRR data is determined until the absolute value of this slope is bellow 0.025, that is the index used for the first point, i.e. the first red star in~Fig.~\ref{Fig:MCC_SI_baseline}. The second red star is identified as the first place where the change in slope is sufficiently small. A 90-point sliding window is moved from the 100th data point until the green dot. For each window the slope of a linear fit of the temperature-HRR is determined. The second red star is the first point for which the change in slope between windows is smaller than 0.0015. 


In the fourth step, the data points used for the baseline correction located ater the HRR peak are determined. The algorithm identifies where the HRR curve becomes essentially stable again, using two independent methods and then taking the more conservative result. For both methods a 50-point window moves between the green point after the peak HRR and the end of the data (excluding the last 100 data points). For each window the slope of a linear fit of Temperature-HRR data is determined. In method one, the difference in slope needs to become sufficiently small, i.e. the first point for which the difference is bellow 0.00005 is selected, unless this does not exist then the difference should be smaller than 0.001. In method 2, the slope of the smoothed (an additional first order savitsky golay filter with 19 points is applied)  data should be larger than -0.01 but smaller than 0.001. The largest index between method 1 and 2 is taken as the baseline point, i.e. as the first blue triangle in Fig.~\ref{Fig:MCC_SI_baseline}. The second blue triangle is 40 data points before the end of the test. It is avoided to take data from the last seconds in the test, since several datasets have a sudden jump near the end. 


Finally, a third degree polynomial is fitted through the regions found in the previous two steps, i.e. all data between the red stars and all data between the blue triangles in Fig.~\ref{Fig:MCC_SI_baseline}. Next, this baseline is subtracted from the HRR data and data outside of the baseline correction is removed. A third degree polynomial was selected because it fitted the data best, considered alternatives include a linear fit and a 1/T fit. 

\section{Cone Calorimeter}
\subsection{Heat Release Rate (HRR)}
\label{Cone_HRR}
\subsubsection{External heat flux: 25 kW/m$^2$}
\begin{minipage}{0.5\textwidth}
\begin{figure}[H]
{\includegraphics[width=3.25in]{../SCRIPT_FIGURES/Cone/Average/IMT_Wood_Cone_25kW_hor.pdf}}\\
\end{figure}
\end{minipage}


\subsubsection{External heat flux: 30 kW/m$^2$}
\begin{minipage}{0.5\textwidth}
\begin{figure}[H]
{\includegraphics[width=3.25in]{../SCRIPT_FIGURES/Cone/Average/TUBS_Wood_Cone_30kW_hor.pdf}}\\
\end{figure}
\end{minipage}
\begin{minipage}{0.35\textwidth}
\begin{figure}[H]
{\includegraphics[width=3.25in]{../SCRIPT_FIGURES/Cone/Average/UDRI_Wood_Cone_30kW_hor.pdf}}\\
\end{figure}
\end{minipage}\\
\begin{minipage}{0.5\textwidth}
\begin{figure}[H]
{\includegraphics[width=3.25in]{../SCRIPT_FIGURES/Cone/Average/UMET_Wood_Cone_30kW_hor.pdf}}\\
\end{figure}
\end{minipage}
\vfill
\newpage
\noindent

\subsubsection{External heat flux: 50 kW/m$^2$}
\begin{minipage}{0.5\textwidth}
\begin{figure}[H]
{\includegraphics[width=3.25in]{../SCRIPT_FIGURES/Cone/Average/IMT_Wood_Cone_50kW_hor.pdf}}\\
\end{figure}
\end{minipage}\begin{minipage}{0.5\textwidth}
\begin{figure}[H]
	{\includegraphics[width=3.25in]{../SCRIPT_FIGURES/Cone/Average/FPL_Wood_Cone_50kW_hor.pdf}}\\
\end{figure}
\end{minipage}
\vfill

\subsubsection{External heat flux: 60 kW/m$^2$}
\begin{minipage}{0.5\textwidth}
\begin{figure}[H]
{\includegraphics[width=3.25in]{../SCRIPT_FIGURES/Cone/Average/UDRI_Wood_Cone_60kW_hor.pdf}}\\
\end{figure}
\end{minipage}
\begin{minipage}{0.35\textwidth}
\begin{figure}[H]
{\includegraphics[width=3.25in]{../SCRIPT_FIGURES/Cone/Average/UMET_Wood_Cone_60kW_hor.pdf}}\\
\end{figure}
\end{minipage}\\
\vfill
\subsubsection{External heat flux: 75 kW/m$^2$}
\begin{minipage}{0.5\textwidth}
\begin{figure}[H]
{\includegraphics[width=3.25in]{../SCRIPT_FIGURES/Cone/Average/IMT_Wood_Cone_75kW_hor.pdf}}\\
\end{figure}
\end{minipage}\\
\vfill 
\newpage

\subsection{Surface Temperature}
\label{Cone_Temp}
Temperatures measured in the same repetition experiment are plotted in the same color. Different line styles are used to indicate measurements recorded at different locations. All plotted temperatures represent backside surface temperature unless otherwise specified in the figure caption.

\subsubsection{External heat flux: 30 kW/m$^2$}
\begin{minipage}{0.5\textwidth}
\begin{figure}[H]
{\includegraphics[width=3.25in]{../SCRIPT_FIGURES/Cone/Individual/Cone_UDRI_Wood_Cone_30kW_hor_BackT.pdf}}\\
\end{figure}
\end{minipage}
\vfill

\subsubsection{External heat flux: 50 kW/m$^2$}
\begin{minipage}{0.5\textwidth}
\begin{figure}[H]
{\includegraphics[width=3.25in]{../SCRIPT_FIGURES/Cone/Individual/Cone_FPL_Wood_Cone_50kW_hor_BackT.pdf}}\\
\caption*{Temperatures were measured at the top surface (dashed line) and at depths 8.6 mm (solid line) and 17.2 mm (dotted line).}
\end{figure}
\end{minipage}
\vfill

\subsubsection{External heat flux: 60 kW/m$^2$}
\begin{minipage}{0.5\textwidth}
	\begin{figure}[H]
		{\includegraphics[width=3.25in]{../SCRIPT_FIGURES/Cone/Individual/Cone_UDRI_Wood_Cone_60kW_hor_BackT.pdf}}\\
	\end{figure}
\end{minipage}
\vfill


\newpage
\section{Anaerobic Gasification}
\subsection{Gasification Mass Flux}
\label{Gas_Mass}
\subsubsection{External heat flux: 30 kW/m$^2$}
\begin{minipage}{0.5\textwidth}
\begin{figure}[H]
{\includegraphics[width=3.25in]{../SCRIPT_FIGURES/Cone/Average/TIFP+UCT_Wood_Gasification_30kW_hor_parallel.pdf}}\\
\end{figure}
\end{minipage}
\begin{minipage}{0.35\textwidth}
\begin{figure}[H]
{\includegraphics[width=3.25in]{../SCRIPT_FIGURES/Cone/Average/TIFP+UCT_Wood_Gasification_30kW_hor_perpendicular.pdf}}\\
\end{figure}
\end{minipage}\\
\vfill

\newpage
\subsubsection{External heat flux: 40 kW/m$^2$}
\begin{minipage}{0.5\textwidth}
\begin{figure}[H]
{\includegraphics[width=3.25in]{../SCRIPT_FIGURES/Cone/Average/FSRI_Wood_CAPA_N2_40kW.pdf}}\\
\end{figure}
\end{minipage}\\
\vfill

\subsubsection{External heat flux: 60 kW/m$^2$}
\begin{minipage}{0.5\textwidth}
\begin{figure}[H]
{\includegraphics[width=3.25in]{../SCRIPT_FIGURES/Cone/Average/TIFP+UCT_Wood_Gasification_60kW_hor_parallel.pdf}}\\
\end{figure}
\end{minipage}
\begin{minipage}{0.35\textwidth}
\begin{figure}[H]
{\includegraphics[width=3.25in]{../SCRIPT_FIGURES/Cone/Average/TIFP+UCT_Wood_Gasification_60kW_hor_perpendicular.pdf}}\\
\end{figure}
\end{minipage}\\
\begin{minipage}{0.5\textwidth}
\begin{figure}[H]
{\includegraphics[width=3.25in]{../SCRIPT_FIGURES/Cone/Average/FSRI_Wood_CAPA_N2_60kW.pdf}}\\
\end{figure}
\end{minipage}\\
\vfill

\newpage
\subsection{Surface Temperature}
\label{Gas_Temp}
Temperatures measured in the same repetition experiment are plotted in the same color with different line styles. All plotted temperatures represent backside surface temperature unless otherwise specified in the figure caption.
\subsubsection{External heat flux: 30 kW/m$^2$}
\begin{minipage}{0.5\textwidth}
\begin{figure}[H]
{\includegraphics[width=3.25in]{../SCRIPT_FIGURES/Cone/Individual/Gasification_TIFP+UCT_Wood_Gasification_30kW_hor_parallel_BackT.pdf}}\\
\caption*{All measurements were made at the backside of the sample, the thermocouples were located along the diagonal of the square samples. Solid lines were measured toward the top-right corner (7.1~cm, 7.1~cm), dotted lines at the center of the square (5.0~cm, 5.0~cm), dashed-dotted lines toward the bottom-left corner (2.9~cm, 2.9~cm).}
\end{figure}
\end{minipage}
\begin{minipage}{0.35\textwidth}
\begin{figure}[H]
{\includegraphics[width=3.25in]{../SCRIPT_FIGURES/Cone/Individual/Gasification_TIFP+UCT_Wood_Gasification_30kW_hor_perpendicular_BackT.pdf}}\\
\caption*{All measurements were made at the backside of the sample, the thermocouples were located along the diagonal of the square samples. Solid lines were measured toward the top-right corner (7.1~cm, 7.1~cm), dotted lines at the center of the square (5.0~cm, 5.0~cm), dashed-dotted lines toward the bottom-left corner (2.9~cm, 2.9~cm).}
\end{figure}
\end{minipage}\\
\vfill


\subsubsection{External heat flux: 40 kW/m$^2$}
\begin{minipage}{0.5\textwidth}
\begin{figure}[H]
{\includegraphics[width=3.25in]{../SCRIPT_FIGURES/Cone/Individual/Gasification_FSRI_Wood_CAPA_N2_40kW_BackT.pdf}}\\
\end{figure}
\end{minipage}\\
Measured front surface (dotted line) and back surface (solid line) temperature during CAPA experiments conducted at 40 kW/m$^2$. 
\vfill
\newpage
\subsubsection{External heat flux: 60 kW/m$^2$}
\begin{minipage}{0.5\textwidth}
\begin{figure}[H]
{\includegraphics[width=3.25in]{../SCRIPT_FIGURES/Cone/Individual/Gasification_TIFP+UCT_Wood_Gasification_60kW_hor_parallel_BackT.pdf}}\\
\caption*{All measurements were made at the backside of the sample, the thermocouples were located along the diagonal of the square samples. Solid lines were measured toward the top-right corner (7.1~cm, 7.1~cm), dotted lines at the center of the square (5.0~cm, 5.0~cm), dashed-dotted lines toward the bottom-left corner (2.9~cm, 2.9~cm).}
\end{figure}
\end{minipage}
\begin{minipage}{0.35\textwidth}
\begin{figure}[H]
{\includegraphics[width=3.25in]{../SCRIPT_FIGURES/Cone/Individual/Gasification_TIFP+UCT_Wood_Gasification_60kW_hor_perpendicular_BackT.pdf}}\\
\caption*{All measurements were made at the backside of the sample, the thermocouples were located along the diagonal of the square samples. Solid lines were measured toward the top-right corner (7.1~cm, 7.1~cm), dotted lines at the center of the square (5.0~cm, 5.0~cm), dashed-dotted lines toward the bottom-left corner (2.9~cm, 2.9~cm).}
\end{figure}
\end{minipage}\\
\begin{minipage}{0.5\textwidth}
\begin{figure}[H]
{\includegraphics[width=3.25in]{../SCRIPT_FIGURES/Cone/Individual/Gasification_FSRI_Wood_CAPA_N2_60kW_BackT.pdf}}\\
\caption*{Measured front surface (dotted line) and back surface (solid line) temperature during CAPA experiments conducted at 60 kW/m$^2$. }
\end{figure}
\end{minipage}\\\\




